\chapter{Dependencias del entorno virtual}
\label{app:anexo_entorno} 

En este anexo se detalla la configuración exacta del entorno de software utilizado para el desarrollo, entrenamiento y despliegue de los modelos de inteligencia artificial del proyecto. La gestión del entorno se ha llevado a cabo mediante el gestor de dependencias \texttt{uv}.

\section*{Entorno base}
\begin{itemize}
    \item \textbf{Lenguaje principal:} Python 3.11.14
    \item \textbf{Gestor de paquetes:} uv
\end{itemize}

\section*{Listado de librerías}
A continuación, se expone el listado exhaustivo de las dependencias principales requeridas para la correcta ejecución del código fuente, tal y como se especificarían en el archivo de configuración del proyecto:

\begin{verbatim}
fastapi>=0.135.1
ipykernel>=7.2.0
joblib>=1.5.3
jupyter>=1.1.1
librosa>=0.11.0
matplotlib>=3.10.8
noisereduce>=3.0.3
numpy>=2.4.2
pandas>=3.0.1
praat-parselmouth>=0.4.7
python-multipart>=0.0.22
scikit-learn>=1.8.0
scipy>=1.17.0
seaborn>=0.13.2
shap>=0.49.1
timm>=1.0.25
torch>=2.6.0
torchvision>=0.20.0
torchaudio>=2.6.0
tqdm>=4.67.3
uvicorn>=0.41.0
xgboost==2.0.3
openpyxl>=3.1.5
transformers>=5.7.0
\end{verbatim}